% !TeX root = ../Chinese Philosophy.tex
\section{中国哲学概述}

本部分主要介绍中国哲学的研究内容和研究主题。

\subsection{中国哲学}

首先是现代视野下中国学者冯友兰的初步认识，分为三个部分：

\begin{itemize}
    \item 宇宙论
    \begin{itemize}
        \item 本体论：存在之本体和真实之要素
        \item 宇宙论：世界发生和历史归宿
    \end{itemize}
    \item 人生论
    \begin{itemize}
        \item 人究竟是什么
        \item 伦理学和政治社会哲学
    \end{itemize}
    \item 知识论
\end{itemize}

张岱年的初步认识有些许不同，但是方向是一样的。冯友兰提出：中国哲学是中国古代的所有学问中，可以被纳入西方对哲学定义中的部分的。随着现代学术的发展和中国哲学知识体系的建构和探索，冯友兰，张岱年和李存山等人发展出了新的中哲内容。

\subsection{中国哲学史}

中国哲学的内容离不开哲学史的积累。冯友兰提出中国哲学的“照着讲”和“接着讲”。所谓接着讲指的是中国哲学的传统接续了宋明以后道学中的理学一派，一度称之为新理学。照着讲更倾向于哲学史的研究，呈现原来的观点。

中国哲学史研究的意义，借助牟宗三的表述和提示：中国思想发展历史以儒家为主干，受到原始模型的笼罩，也引申为原始模型的规范。汉代继承，魏晋道家复兴，接续佛学传入，从而实现充实和丰富。（所谓歧出）宋明归于自己，继承儒家发展进入新阶段，直至今日和西方哲学重新碰撞和丰富。

中国哲学史的研究方法：\textbf{始终存在的方法问题}。冯友兰指出中哲略于方法组织，近人多诟病，然而这是中哲的精神。他认为中国哲学虽然没有形式上的系统但是有实质上的系统（中国哲学没有不成东西），与西方哲学是共有的。

同时，研究中国哲学史必须加以解释和排比，在解释和整理的时候往往在表达自己的观点和对于古人的揣测。中国哲学的研究任务中就要努力找知识谱系（研究梳理中国古代哲学史，对古人要具备了解之同情），并努力排除个人的成见，经验和习惯（所谓神游冥想，处于统一境界），理解其用意（之所以不得不如是之苦心孤诣），了解并批评学说的是非。

\subsection{分期和阶段}

\textbf{分期：}先秦诸子学，两汉经学，魏晋玄学，南北朝隋唐佛学，宋明理学，清代和近代哲学（先秦到魏晋为中国哲学上）

